problem:  
Let $\Pi_1,\Pi_2,\Pi_3\subset \mathbb{R}^3$ be three half-planes bounded by a common line $\ell$, meeting symmetrically at $120^\circ$ along $\ell$. Denote by $\mathfrak{M}_{\text{triple}}$ the moduli space (possibly empty) of properly embedded minimal surfaces $M\subset \mathbb{R}^3$ such that:  

- $M$ is asymptotic to $\Pi_1\cup \Pi_2\cup \Pi_3$ near infinity;  
- $M$ is smooth and embedded away from $\ell$, and the asymptotic intersection of its three ends with $\ell$ reproduces the $120^\circ$ configuration.  

**Question:**  
1. Does there exist *any* such minimal surface $M\in \mathfrak{M}_{\text{triple}}$?  
2. If so, is the moduli space $\mathfrak{M}_{\text{triple}}$ finite-dimensional and smooth near such $M$, and what is the expected dimension of its Jacobi field space modulo Euclidean isometries?  
3. Can such a local model (if it exists) be geometrically and analytically adapted—via scaling charts and matching of normal bundles—to desingularize a triple junction configuration of compact minimal surfaces in $S^3$ meeting along a geodesic with equal $120^\circ$ angles?  

In particular, is it theoretically possible to carry out a gluing construction in $S^3$ replacing a small tubular neighborhood of the triple intersection region by a rescaled copy of a local model $M\in \mathfrak{M}_{\text{triple}}$ (should it exist), thereby producing a smooth embedded minimal surface $\Sigma_\varepsilon$ in $S^3$ that approximates the union $\Sigma_1\cup\Sigma_2\cup\Sigma_3$ as $\varepsilon\to 0$?  

Why is it a "good" problem:  
This refined formulation isolates the core, unresolved question of whether a **local minimal surface model for a $120^\circ$ triple junction** exists—and if it does, whether it can serve as a building block in global minimal surface constructions in $S^3$.  

1. **Profound Insight:**  
It distills the essential analytic unknown in the extension of desingularization theory from two-sheet intersections to the *first genuinely new local singular type* beyond orthogonal crossings. The existence (or nonexistence) of such a local minimal model would reveal whether the minimal surface equation admits smoothing of symmetric triple-junction singularities, a question fundamental to the structure of the minimal surface landscape.  

2. **Bridge Across Fields:**  
The problem connects nonlinear elliptic PDEs (local minimal surface equations and Jacobi field analysis), geometric measure theory (limits of soap-film-type minimal networks), and global differential geometry (applications in the three-sphere). In doing so, it creates a conceptual bridge between **classical minimal surface gluing** and **Plateau-type singular network geometries**.  

3. **Extreme Simplicity:**  
The statement is geometrically vivid and succinct: “Does a minimal surface exist asymptotic to three half-planes meeting at $120^\circ$?” Conceptually simple, yet it addresses an uncharted region in analytical and geometric theory.  

4. **Potential Impact:**  
If an embedded minimal surface in $\mathbb{R}^3$ asymptotic to a symmetric $120^\circ$ triple junction exists, it would introduce a *new local model*—analogous to Scherk surfaces for two-sheet intersections. Such a discovery would revolutionize the gluing approach in $S^3$, allowing multi-surface desingularizations and advancing our understanding of how minimal surfaces can interact and “fuse.” Conversely, proving nonexistence would uncover a fundamental rigidity barrier in the geometry of minimal submanifolds, reshaping conjectures on the analytic limits of desingularization.  

Thus, the problem is *good* because it identifies the precise unknown that must be resolved before any triple-junction gluing can exist: the local analytic existence of a smooth minimal model asymptotic to a $120^\circ$ triple-plane configuration. Its formulation is minimalistic, logically clean, and rooted at the frontier between geometric analysis and topology.